\documentclass[11pt,a4paper]{article}

\usepackage[margin=2.3cm]{geometry}
\usepackage[T1]{fontenc}
\usepackage[utf8]{inputenc}
\usepackage{lmodern}
\usepackage{microtype}
\usepackage{enumitem}
\usepackage{booktabs}
\usepackage{tabularx}
\usepackage{longtable}
\usepackage{xcolor}
\usepackage{listings}
\usepackage{hyperref}
\usepackage{fancyhdr}

\definecolor{codeblue}{RGB}{36,86,148}
\definecolor{codegreen}{RGB}{40,120,70}
\definecolor{codegray}{RGB}{105,105,105}
\definecolor{codebackground}{RGB}{247,247,247}

\hypersetup{
  colorlinks=true,
  linkcolor=codeblue,
  urlcolor=codeblue,
  citecolor=codeblue,
  pdftitle={PaiSmart Application and Test Automation Report}
}

\lstset{
  backgroundcolor=\color{codebackground},
  basicstyle=\ttfamily\small,
  breaklines=true,
  columns=fullflexible,
  commentstyle=\color{codegreen},
  frame=single,
  keywordstyle=\color{codeblue}\bfseries,
  numbers=left,
  numbersep=8pt,
  numberstyle=\tiny\color{codegray},
  showstringspaces=false,
  tabsize=2
}

\pagestyle{fancy}
\fancyhf{}
\lhead{PaiSmart Test Automation}
\rhead{Student ID: S2025XXXX}
\cfoot{\thepage}
\setlength{\headheight}{14pt}

\title{\textbf{Application Development and Test Automation Report}\\
\large A Document Knowledge Base Module Based on PaiSmart}
\author{Student Name: XXX\\Student ID: S2025XXXX}
\date{June 7, 2026}

\begin{document}

\maketitle
\thispagestyle{empty}

\begin{abstract}
This report presents a small web application module developed and studied within the
PaiSmart project: a document knowledge base that supports file upload, document
processing, text chunking, indexing, and hybrid retrieval. The report explains how the
module is assembled from a Vue 3 front end and a Spring Boot back end, how its major
components cooperate, and how it can be demonstrated as a self-contained application.
It then describes the design and execution of automated tests for user registration,
authentication, conversation persistence, and text chunking. The executable unit-test
suite uses JUnit 5, Mockito, and Maven Surefire. On June 7, 2026, the selected suite
executed 10 tests with zero failures, zero errors, and zero skipped tests. The report
also evaluates additional automation tools that can extend the project from unit
testing to API, browser, integration, coverage, and performance testing.
\end{abstract}

\tableofcontents
\newpage

\section{Assignment Interpretation and Project Plan}

The assignment requires three connected deliverables:
\begin{enumerate}[leftmargin=2em]
  \item create or identify a small website, application, or software module;
  \item write automated test scripts for important behavior; and
  \item submit a report describing both the application and the testing work.
\end{enumerate}

Instead of claiming that the complete PaiSmart platform was created from the beginning
for this assignment, this work selects and packages one coherent feature as the small
application: the \textbf{Document Knowledge Base and Retrieval Module}. This scope is
large enough to demonstrate full-stack development, but focused enough to explain and
test clearly.

The practical completion plan is:
\begin{enumerate}[leftmargin=2em]
  \item inspect the existing PaiSmart architecture and choose a demonstrable module;
  \item identify the front-end pages, REST controllers, services, repositories, and
        infrastructure used by that module;
  \item prepare a repeatable demonstration of document upload and retrieval;
  \item treat the Java files under
        \path{src/test/java/com/yizhaoqi/smartpai/service/} as automated test scripts;
  \item repair outdated mocks so the tests match the current production code;
  \item run a lightweight unit-test suite that does not require external services;
  \item record actual Maven Surefire results; and
  \item document future API, browser, integration, coverage, and performance tests.
\end{enumerate}

\section{Application: Document Knowledge Base and Retrieval Module}

\subsection{Purpose and User Problem}

The selected application helps a user convert ordinary documents into a searchable
knowledge base. A user can upload a document, monitor upload progress, associate it
with an organization tag, choose whether it is public or private, and search the
processed content. This solves a common problem: information may exist in files, but it
is difficult to find relevant passages quickly by filename alone.

The module provides the following user-visible functions:
\begin{itemize}[leftmargin=2em]
  \item chunked file upload with progress tracking and interrupted-upload recovery;
  \item file type validation and MD5-based file identification;
  \item document listing, preview, deletion, and organization-based access control;
  \item asynchronous parsing and text chunking;
  \item vector and keyword indexing in Elasticsearch; and
  \item hybrid knowledge-base search with ranked passages and source filenames.
\end{itemize}

\subsection{Technology Stack}

\begin{table}[htbp]
\centering
\caption{Technology stack used by the selected module}
\begin{tabularx}{\textwidth}{@{}lX@{}}
\toprule
\textbf{Layer} & \textbf{Technology and responsibility} \\
\midrule
Front end &
Vue 3.5, TypeScript, Vite 6, Naive UI, Pinia, Vue Router, and pnpm. The front end
implements the knowledge-base page, upload dialog, search dialog, progress display,
and API calls. \\
Back end &
Java 17 and Spring Boot 3.4.2. Spring MVC exposes REST endpoints; service classes
implement business logic; Spring Data JPA accesses relational data. \\
Database &
MySQL 8 stores users, uploaded-file metadata, conversations, organization tags, and
document-related records. \\
Cache and upload state &
Redis stores temporary upload state, including bitmap information for uploaded chunks. \\
Object storage &
MinIO stores file chunks and merged documents. \\
Asynchronous processing &
Apache Kafka transfers file-processing tasks after chunks have been merged. \\
Document processing &
Apache Tika extracts text; HanLP supports Chinese-language text segmentation. \\
Search &
Elasticsearch 8.10 stores indexed document content and supports retrieval. \\
Security &
Spring Security and JWT protect endpoints and identify the current user. \\
Testing &
JUnit 5, Mockito, Spring Boot Test, Maven Surefire, and Spring test utilities. \\
\bottomrule
\end{tabularx}
\end{table}

\subsection{Application Architecture}

The module follows a layered architecture:
\begin{enumerate}[leftmargin=2em]
  \item \textbf{Presentation layer.} Vue components display the document table and
        dialogs. The key files are
        \path{frontend/src/views/knowledge-base/index.vue},
        \path{modules/upload-dialog.vue}, and
        \path{modules/search-dialog.vue}.
  \item \textbf{API layer.} Controllers such as \texttt{UploadController},
        \texttt{DocumentController}, and \texttt{SearchController} receive HTTP
        requests, validate request data, and return structured responses.
  \item \textbf{Service layer.} \texttt{UploadService}, \texttt{ParseService},
        \texttt{VectorizationService}, \texttt{HybridSearchService}, and
        \texttt{DocumentService} implement the workflow.
  \item \textbf{Persistence layer.} Spring Data repositories store file metadata,
        document vectors, chunks, users, and conversations.
  \item \textbf{Infrastructure layer.} MySQL, Redis, MinIO, Kafka, Elasticsearch, and
        external embedding services provide storage and processing capabilities.
\end{enumerate}

This separation makes individual business services testable with mocks. For example,
\texttt{ConversationService} can be tested without starting MySQL because
\texttt{UserRepository} and \texttt{ConversationRepository} are replaced by Mockito
objects.

\subsection{How the Small Application Was Created}

The small application is produced by selecting the knowledge-base feature from the
larger PaiSmart system and organizing its implementation into a complete user journey.
The development process is as follows.

\subsubsection{Step 1: Define the functional scope}

The scope includes document upload, processing, management, and retrieval. Chat,
administration, and unrelated user-interface pages remain part of PaiSmart but are not
the main demonstration target.

\subsubsection{Step 2: Build the front-end workflow}

The knowledge-base page loads document records and displays filename, MD5, size,
upload status, organization tag, visibility, and creation time. The upload dialog
collects a file, organization tag, and public/private setting. The front end calculates
an MD5 value, divides a large file into chunks, sends chunks to the server, and can
resume an interrupted upload. The search dialog sends a query and a \texttt{topK}
parameter to the hybrid-search endpoint, then displays ranked text passages, scores,
and source filenames.

\subsubsection{Step 3: Implement the upload API}

\texttt{UploadController} exposes endpoints under \path{/api/v1/upload}. The chunk
endpoint receives file metadata and a multipart chunk. The controller validates the
file type, determines the user's organization tag when necessary, delegates storage to
\texttt{UploadService}, and reports upload progress.

After all chunks arrive, the merge endpoint:
\begin{enumerate}[leftmargin=2em]
  \item checks that the file belongs to the current user;
  \item confirms that all expected chunks have been uploaded;
  \item asks \texttt{UploadService} to compose the chunks in MinIO;
  \item creates a \texttt{FileProcessingTask}; and
  \item publishes the task to Kafka for asynchronous parsing and indexing.
\end{enumerate}

\subsubsection{Step 4: Parse and index the document}

The file-processing consumer retrieves the merged document. \texttt{ParseService}
extracts text and splits long content into manageable chunks. Each text chunk can then
be transformed into an embedding by \texttt{VectorizationService}. Searchable content
and vector information are stored for later retrieval.

\subsubsection{Step 5: Implement retrieval}

The search page calls the hybrid-search endpoint. The back end combines retrieval
signals, ranks matching chunks, applies access-control rules, and returns results. The
front end highlights matching query text and shows the source document for traceability.

\subsubsection{Step 6: Add security and multi-tenant separation}

JWT authentication identifies the caller. Organization tags and visibility fields
restrict which documents a user can access. A document can be private to an
organization or public according to its metadata.

\subsection{End-to-End Execution Flow}

The application flow can be summarized as:
\begin{center}
\texttt{Select file}
$\rightarrow$ \texttt{calculate MD5}
$\rightarrow$ \texttt{upload chunks}
$\rightarrow$ \texttt{merge in MinIO}
$\rightarrow$ \texttt{publish Kafka task}
$\rightarrow$ \texttt{parse text}
$\rightarrow$ \texttt{split and vectorize}
$\rightarrow$ \texttt{index}
$\rightarrow$ \texttt{hybrid search}.
\end{center}

\subsection{Running and Demonstrating the Application}

The infrastructure can be started from the supplied Docker Compose configuration:
\begin{lstlisting}[language=bash]
docker compose -f docs/docker-compose.yaml up -d
\end{lstlisting}

The back end can then be started from the project root:
\begin{lstlisting}[language=bash]
mvn spring-boot:run
\end{lstlisting}

The front end can be started separately:
\begin{lstlisting}[language=bash]
cd frontend
pnpm install
pnpm dev
\end{lstlisting}

A suitable demonstration is:
\begin{enumerate}[leftmargin=2em]
  \item sign in to the application;
  \item open the knowledge-base page;
  \item upload a small PDF, DOCX, or TXT document;
  \item observe upload and processing status;
  \item open the search dialog;
  \item enter a phrase that appears in the document; and
  \item show the returned passage, relevance score, and source filename.
\end{enumerate}

\section{Automated Test Report}

\subsection{Test Objectives}

The automated tests are intended to verify deterministic business rules quickly and
repeatably. The selected unit-test suite checks:
\begin{itemize}[leftmargin=2em]
  \item successful user registration;
  \item rejection of a duplicate username;
  \item successful password authentication;
  \item rejection of invalid credentials;
  \item persistence of a conversation record;
  \item retrieval of a user's conversation history;
  \item preservation of text when long sentences are divided into chunks;
  \item empty, short, and long text boundary behavior;
  \item compliance with different chunk-size limits; and
  \item basic execution-time behavior for text splitting.
\end{itemize}

\subsection{Existing Automated Test Scripts}

The service test directory contains five relevant test classes:

\begin{longtable}{@{}p{0.29\textwidth}p{0.18\textwidth}p{0.45\textwidth}@{}}
\caption{Classification of existing service tests}\\
\toprule
\textbf{Test class} & \textbf{Type} & \textbf{Purpose and execution conditions} \\
\midrule
\endfirsthead
\toprule
\textbf{Test class} & \textbf{Type} & \textbf{Purpose and execution conditions} \\
\midrule
\endhead
\texttt{ConversationServiceTest} &
Unit test &
Uses Mockito repositories to verify conversation saving and retrieval without starting
Spring or a database. \\
\texttt{UserServiceTest} &
Unit test &
Uses Mockito to isolate registration and authentication logic. It verifies repository
interactions, password hashing, roles, private organization tags, cache updates, and
error responses. \\
\texttt{ParseServiceUnitTest} &
Unit test &
Creates \texttt{ParseService} directly and invokes the private splitting method through
reflection. It checks reconstruction, edge cases, size limits, and a broad timing
threshold. \\
\texttt{ParseServiceTest} &
Spring-context test &
Uses \texttt{@SpringBootTest}. It exercises similar splitting behavior but attempts to
load the application context, so it is heavier and may depend on configuration. \\
\texttt{UploadServicePerformanceTest} &
Illustrative Spring test &
Logs a calculated comparison between repeated Redis calls and one bitmap retrieval.
It is useful documentation, but it is not a real benchmark because it does not measure
the production methods against a live or simulated Redis workload. \\
\bottomrule
\end{longtable}

The first three classes form the repeatable assignment suite because they execute
without MySQL, Redis, Kafka, MinIO, Elasticsearch, or a running Spring application.

\subsection{How the Test Automation Scripts Were Written}

\subsubsection{Arrange--Act--Assert structure}

Each test follows the Arrange--Act--Assert pattern:
\begin{itemize}[leftmargin=2em]
  \item \textbf{Arrange:} create model objects and define mock responses;
  \item \textbf{Act:} call one public service method; and
  \item \textbf{Assert:} check returned values, exceptions, state, or dependency calls.
\end{itemize}

\subsubsection{Mocking external dependencies}

Mockito annotations isolate the class under test:
\begin{lstlisting}[language=Java]
@Mock
private ConversationRepository conversationRepository;

@Mock
private UserRepository userRepository;

@InjectMocks
private ConversationService conversationService;
\end{lstlisting}

The test defines repository behavior and then verifies the service interaction:
\begin{lstlisting}[language=Java]
when(userRepository.findByUsername("testuser"))
    .thenReturn(Optional.of(user));

conversationService.recordConversation(
    "testuser",
    "What is AI?",
    "AI stands for Artificial Intelligence.");

verify(conversationRepository, times(1))
    .save(any(Conversation.class));
\end{lstlisting}

\subsubsection{Testing normal and exceptional paths}

The user-service tests cover both successful and unsuccessful behavior. For duplicate
registration, the repository returns an existing user and the test asserts both the
exception message and HTTP status:
\begin{lstlisting}[language=Java]
CustomException exception = assertThrows(
    CustomException.class,
    () -> userService.registerUser("testuser", "password123"));

assertEquals("Username already exists", exception.getMessage());
assertEquals(HttpStatus.BAD_REQUEST, exception.getStatus());
\end{lstlisting}

\subsubsection{Capturing saved entities}

Registration has no return value, so an \texttt{ArgumentCaptor} inspects the object sent
to the repository. The test verifies that the password is not stored as plain text,
that it matches through \texttt{PasswordUtil}, and that role and organization fields
are assigned correctly.

\subsubsection{Testing text chunking}

\texttt{ParseServiceUnitTest} invokes \texttt{splitLongSentence} through reflection
because the method is private. The most important invariant is lossless reconstruction:
\begin{lstlisting}[language=Java]
List<String> result =
    (List<String>) method.invoke(parseService, sentence, chunkSize);

String reconstructed = String.join("", result);
assertEquals(sentence, reconstructed);
\end{lstlisting}

Although reflection is acceptable for the current assignment, a future refactoring
could move chunking into a dedicated package-private component. That would make the
behavior directly testable without exposing it as a public API.

\subsection{Repair Made to the Existing Test Suite}

The current \texttt{UserService.registerUser} method depends on
\texttt{OrganizationTagRepository} and \texttt{OrgTagCacheService}. The previous test
mocked only \texttt{UserRepository}; therefore, the successful registration test threw
a \texttt{NullPointerException}. The automated script was updated to:
\begin{itemize}[leftmargin=2em]
  \item mock \texttt{OrganizationTagRepository};
  \item mock \texttt{OrgTagCacheService};
  \item define behavior for the default and private organization tags;
  \item expect two user-save operations performed by the current implementation;
  \item verify password hashing, role, private tag, and primary organization; and
  \item verify organization-tag creation and cache updates.
\end{itemize}

This repair is important because an automated test is useful only when it reflects the
current production contract.

\subsection{Test Cases and Expected Results}

\begin{longtable}{@{}p{0.08\textwidth}p{0.30\textwidth}p{0.50\textwidth}@{}}
\caption{Executed automated test cases}\\
\toprule
\textbf{ID} & \textbf{Test method} & \textbf{Expected result} \\
\midrule
\endfirsthead
\toprule
\textbf{ID} & \textbf{Test method} & \textbf{Expected result} \\
\midrule
\endhead
UT-01 & \texttt{testRegisterUser\_Success} &
The user is saved with a hashed password, USER role, private organization tag, and
cache entries. \\
UT-02 & \texttt{testRegisterUser\_UsernameExists} &
A duplicate username produces \texttt{400 BAD REQUEST}. \\
UT-03 & \texttt{testAuthenticateUser\_Success} &
The correct raw password matches the stored hash and returns the username. \\
UT-04 & \texttt{testAuthenticateUser\_InvalidCredentials} &
An unknown user produces \texttt{401 UNAUTHORIZED}. \\
UT-05 & \texttt{testRecordConversation} &
One conversation entity is passed to the repository's save method. \\
UT-06 & \texttt{testGetConversations} &
The service returns the repository's conversation list for the user ID. \\
UT-07 & \texttt{testSplitLongSentence\_BasicFunctionality} &
The result is nonempty and reconstructs the original sentence. \\
UT-08 & \texttt{testSplitLongSentence\_EdgeCases} &
Empty, one-character, and long inputs are handled without content loss. \\
UT-09 & \texttt{testSplitLongSentence\_ChunkSizeValidation} &
Chunks respect several requested sizes and reconstruct the source text. \\
UT-10 & \texttt{testSplitLongSentence\_Performance} &
A 2,900-character input is divided successfully in less than five seconds. \\
\bottomrule
\end{longtable}

\subsection{Test Execution Command}

The verified lightweight suite was executed with Java 17:
\begin{lstlisting}[language=bash]
mvn "-Dtest=ConversationServiceTest,UserServiceTest,ParseServiceUnitTest" test
\end{lstlisting}

Individual classes or methods can also be executed:
\begin{lstlisting}[language=bash]
mvn -Dtest=ConversationServiceTest test
mvn -Dtest=UserServiceTest test
mvn -Dtest=ParseServiceUnitTest test
mvn -Dtest=ConversationServiceTest#testRecordConversation test
\end{lstlisting}

\subsection{Verified Test Results}

The suite was executed on \textbf{June 7, 2026 at 12:11 China Standard Time
(UTC+8)}. Maven Surefire reported:

\begin{table}[htbp]
\centering
\caption{Verified Maven Surefire results}
\begin{tabular}{@{}lrrrr@{}}
\toprule
\textbf{Test class} & \textbf{Run} & \textbf{Failures} & \textbf{Errors} & \textbf{Skipped} \\
\midrule
\texttt{ConversationServiceTest} & 2 & 0 & 0 & 0 \\
\texttt{ParseServiceUnitTest} & 4 & 0 & 0 & 0 \\
\texttt{UserServiceTest} & 4 & 0 & 0 & 0 \\
\midrule
\textbf{Total} & \textbf{10} & \textbf{0} & \textbf{0} & \textbf{0} \\
\bottomrule
\end{tabular}
\end{table}

The Maven result was \texttt{BUILD SUCCESS}. Surefire generated machine-readable XML
and text reports under \path{target/surefire-reports/}. These files can be archived by
a continuous-integration server.

\subsection{Analysis of Test Quality and Limitations}

The executed tests are fast and isolated, making them suitable for every local build
and every commit. They verify important service behavior without failures caused by
unavailable infrastructure.

However, the current coverage is not complete:
\begin{itemize}[leftmargin=2em]
  \item the upload controller and upload service are not covered by realistic unit tests;
  \item no test verifies Kafka publication after a successful merge;
  \item no integration test verifies MinIO composition, Redis bitmaps, or Elasticsearch;
  \item no browser test performs the complete upload-and-search workflow;
  \item the conversation tests do not yet cover missing users, date filters, or
        administrator authorization;
  \item the authentication tests should add an existing user with a wrong password; and
  \item the current performance test uses a generous threshold and is not a statistically
        reliable benchmark.
\end{itemize}

\section{At Least Five Available Test Automation Tools and Their Functions}

This section lists at least five available test automation tools and discusses their
main functions and possible applications in PaiSmart. Eight tools are included to
cover unit testing, test execution, code coverage, API testing, end-to-end testing,
integration testing, and performance testing.

\begin{longtable}{@{}p{0.20\textwidth}p{0.28\textwidth}p{0.44\textwidth}@{}}
\caption{Recommended automation tools}\\
\toprule
\textbf{Tool} & \textbf{Testing level} & \textbf{Function in PaiSmart} \\
\midrule
\endfirsthead
\toprule
\textbf{Tool} & \textbf{Testing level} & \textbf{Function in PaiSmart} \\
\midrule
\endhead
JUnit 5 &
Unit and integration testing &
Defines Java test cases, lifecycle methods, assertions, parameterized tests, tags, and
test suites. It is the primary framework used by the existing service tests. \\
Mockito &
Unit testing &
Creates mock repositories and services, configures return values, captures arguments,
and verifies interactions. It keeps business tests independent of MySQL, Redis, and
other systems. \\
Maven Surefire &
Test execution and reporting &
Discovers JUnit tests during the Maven \texttt{test} phase and produces console, text,
and XML results suitable for CI pipelines. \\
JaCoCo &
Code coverage &
Measures line and branch coverage. It can reveal untested paths in upload, parsing,
security, and search services and can enforce a minimum coverage threshold. \\
Postman and Newman &
REST API testing &
Postman defines collections for login, upload status, merge, document listing, and
search endpoints. Newman runs the same collection automatically from the command line
or CI. \\
Playwright &
End-to-end browser testing &
Automates Chromium, Firefox, or WebKit to test login, file selection, upload progress,
search, preview, and deletion from the user's point of view. \\
Testcontainers &
Integration testing &
Starts disposable MySQL, Redis, Kafka, MinIO, or Elasticsearch containers for tests,
providing realistic infrastructure without relying on manually maintained shared
services. \\
JMeter &
Load and performance testing &
Generates concurrent upload and search traffic, measures throughput and latency, and
helps evaluate Redis bitmap optimization and Elasticsearch query performance. \\
\bottomrule
\end{longtable}

\section{Recommended Next Automation Stages}

\subsection{Stage 1: Improve Unit Coverage}

Add tests for missing users, wrong passwords, administrator access, date-filtered
conversation queries, file type validation, upload progress calculation, and search
result mapping. These tests should remain independent of external services.

\subsection{Stage 2: Add API Automation}

Create a Postman collection with environment variables for the base URL and JWT token.
The collection should automate:
\begin{enumerate}[leftmargin=2em]
  \item registration and login;
  \item supported-file-type retrieval;
  \item upload-status queries;
  \item document listing and deletion;
  \item hybrid search; and
  \item authorization failures for inaccessible documents.
\end{enumerate}

Newman can execute the collection in CI and generate an HTML or JUnit-style report.

\subsection{Stage 3: Add Containerized Integration Tests}

Use Testcontainers to start infrastructure for tests. A focused integration test can
upload chunks, verify Redis state, merge content in MinIO, consume a Kafka task, index
document content, and query Elasticsearch. These tests should run separately from the
fast unit tests, for example through Maven Failsafe.

\subsection{Stage 4: Add Browser End-to-End Tests}

Use Playwright to automate the same flow used in the classroom demonstration. Stable
\texttt{data-testid} attributes should be added to important controls so tests do not
depend on translated labels or CSS structure.

\subsection{Stage 5: Add Coverage and Continuous Integration}

Configure JaCoCo and run the unit suite for every push. A CI pipeline should:
\begin{enumerate}[leftmargin=2em]
  \item compile the Java application;
  \item execute unit tests;
  \item publish Surefire reports;
  \item generate JaCoCo coverage;
  \item build and type-check the Vue application; and
  \item run integration or browser tests in a later job.
\end{enumerate}

\section{Submission Checklist}

Before submission:
\begin{enumerate}[leftmargin=2em]
  \item replace \texttt{XXX} with the real student name;
  \item replace every \texttt{S2025XXXX} with the real student ID;
  \item rename this file to \texttt{studentID\_TestAutomation.tex}, for example
        \texttt{S202512345\_TestAutomation.tex};
  \item run the verified Maven unit-test command again;
  \item compile the LaTeX file and inspect the PDF;
  \item optionally add one or two screenshots of the knowledge-base page and test output;
  \item ensure the \texttt{.tex} file is placed inside the required SQAM project; and
  \item submit before June 12, 2026.
\end{enumerate}

\section{Conclusion}

This work presents the PaiSmart document knowledge base as a focused small web
application. The module demonstrates a realistic full-stack workflow involving Vue,
Spring Boot, MySQL, Redis, MinIO, Kafka, document parsing, vectorization, and
Elasticsearch retrieval. The automated testing work uses the project's existing
service tests as executable scripts and updates the user-service test to match the
current production dependencies.

The verified lightweight suite contains 10 tests and completes successfully without
external infrastructure. This gives the project a practical regression baseline.
JUnit 5, Mockito, and Maven Surefire already support repeatable unit automation, while
JaCoCo, Postman/Newman, Playwright, Testcontainers, and JMeter provide a clear route to
broader coverage. The resulting report therefore documents both a working application
module and a realistic, staged test automation strategy.

\end{document}
